%%=============================================================================
%% 第 1 章 引言
%%=============================================================================

\xmuchapter{引言}{Introduction}
\label{chap:introduction}

\xmusection{研究背景}{Research Background}
\label{sec:background}

脑机接口（Brain-Computer Interface, BCI）技术通过采集和解码脑电信号（Electroencephalogram, EEG），在大脑与外部环境之间建立通信通道，为运动功能障碍患者提供了重建运动功能的可能性\cite{tangermann2012review}。自 1970 年代 Vidal 首次提出 BCI 概念以来\cite{vidal1973toward}，经过半个世纪的发展，BCI 技术已从实验室研究逐步走向临床应用。根据信号诱发机制的不同，BCI 可分为主动式（如运动想象）、反应式（如 P300）和稳态诱发（如 SSVEP）三大范式\cite{lotte2018review}。

运动想象（Motor Imagery, MI）是 BCI 的重要应用方向，用户通过想象肢体运动在大脑感觉运动皮层产生事件相关去同步化（Event-Related Desynchronization, ERD）现象，BCI 系统通过分类不同运动想象任务诱发的 ERD 模式识别用户意图\cite{pfurtscheller1999event}。ERD 现象由 Pfurtscheller 等于 1970 年代末首次发现并系统描述，其定义为特定频段（主要是$\mu$节律 8--13 Hz 和$\beta$节律 13--30 Hz）的功率谱密度在运动准备和执行过程中的抑制性变化\cite{pfurtscheller1999event}。

与其他 BCI 范式相比，运动想象 BCI 具有无需外部刺激、用户可自主控制、信息传输率高等优势\cite{lotte2018review}。P300 范式依赖奇偶刺激呈现，用户注意力易疲劳；SSVEP 范式虽信息传输率高，但需持续视觉刺激，可能引发光敏性癫痫。运动想象 BCI 则通过内源性认知任务实现控制，具有更好的用户体验和长期可用性\cite{lotte2018review}。

然而，运动想象 BCI 的分类性能受多种因素影响，其中时间窗选择是关键因素之一\cite{yang2015time}。BCI Competition 系列竞赛（2003--2015）的数据集为时间窗研究提供了标准化基准，其中 BCI Competition IV 2a 数据集包含九名健康被试在四种运动想象任务（左手、右手、双脚、舌头）下的 EEG 记录，已成为时间窗优化方法的标准评测平台\cite{tangermann2012review}。

在运动想象 BCI 系统中，时间窗选择的重要性源于三方面：首先，ERD 动态特性存在显著的个体间差异，固定时间窗策略难以适应所有被试\cite{zhang2019temporally}；其次，EEG 信号具有非平稳性和低信噪比特性，时间窗长度需要在信息完整性与噪声抑制之间取得平衡\cite{blankertz2008optimizing}；最后，在线 BCI 系统需要在毫秒级延迟内完成决策，时间窗的起始点和终点直接影响系统响应性\cite{tangermann2012review}。因此，自适应时间窗优化是提升运动想象 BCI 分类性能的关键途径。

\xmusection{问题提出}{Problem Statement}
\label{sec:problem_statement}

BCI 技术在实际应用中面临多重约束。首先是\textbf{数据稀缺性}：EEG 数据采集成本高昂，BCI 研究通常仅能获取数十至数百个试次的训练数据，难以支撑深度神经网络等数据密集型模型的有效训练\cite{lotte2007review}。此外，EEG 信号具有显著的\textbf{个体间差异性}（inter-subject variability），要求 BCI 模型必须能够根据特定用户的早期校准数据（通常仅 50--100 个试次）进行快速独立训练\cite{lotte2011regularizing}。

其次是\textbf{设备资源约束}：便携式 BCI 设备通常采用低功耗嵌入式处理器，配备有限的内存（KB 至 MB 级）和存储空间，且依赖电池供电。在线 BCI 系统需要在毫秒级延迟内完成时间窗选择、特征提取和分类预测，复杂的深度学习模型可能引入不可接受的延迟。

时间窗优化方法的发展历程可追溯至 BCI 研究的早期阶段。根据优化策略的不同，现有方法可分为以下四类：

（1）\textbf{固定时间窗方法}：早期 BCI 系统采用基于先验知识的固定时间窗策略，通常选择刺激呈现后 0.5--2.5 s 作为分析区间\cite{mcfarland1997spatial}。该方法实现简单、计算开销低，但无法适应 ERD 动态特性的个体间差异，分类性能受限\cite{zhang2019temporally}。

（2）\textbf{滑动窗方法}：通过在整个时间轴上滑动固定长度的窗口并选择分类性能最优的位置，滑动窗方法在一定程度上提升了时间窗选择的自适应性\cite{yang2015time}。然而，该方法计算复杂度随窗口数量和步长呈线性增长，且难以同时优化时间窗的起始点和终点两个参数。

（3）\textbf{贝叶斯优化方法}：贝叶斯优化通过构建代理模型（如高斯过程）和目标采集函数，在参数空间中进行高效搜索\cite{yang2015time}。该方法在低维参数空间表现良好，但计算复杂度随参数维度增加呈指数增长，难以处理高维优化问题。

（4）\textbf{深度强化学习方法}：近年来，深度强化学习（如 DQN）被引入 BCI 参数优化任务\cite{chen2019dynamic, zhang2021mars}。该方法通过深度神经网络近似状态 - 动作价值函数，理论上可处理任意规模的状态空间。然而，DQN 存在模型复杂度高（百万级参数）、训练不稳定（对超参数敏感）、可解释性差（黑盒决策）等问题\cite{mnih2015human}，与 BCI 系统的数据稀缺性和设备资源约束形成尖锐矛盾。

上述分析表明，现有时间窗优化方法难以同时满足 BCI 系统的轻量化、稳定性和可解释性要求。尽管深度强化学习在 BCI 参数优化中展现出潜力，但其对大量训练数据和高计算资源的依赖，与 BCI 系统中数据稀缺、设备资源受限的现实形成矛盾。因此，如何在保证性能的前提下实现轻量化、可解释的时间窗优化，仍是当前研究的难点。

针对上述问题，本研究的核心研究问题是：\textbf{在运动想象时间窗优化任务中，是否存在一种方法，能够在保持与深度强化学习相当性能的同时，显著降低模型复杂度、提升训练稳定性和可解释性？}

\xmusection{研究内容与意义}{Research Content and Significance}
\label{sec:research_content}

\xmusubsection{研究内容}{Research Content}
\label{subsec:research_content}

针对上述问题，本研究提出一种基于\textbf{表格型 Q-Learning}的运动想象脑电信号时间窗优化方法。核心主张是：在运动想象时间窗优化这一状态空间有限（136 个状态）的特定任务中，表格型 Q-Learning 可在保持与 DQN 相当性能的同时，显著降低模型复杂度，提升系统稳定性和可解释性。

本研究的主要内容包括：

（1）\textbf{问题形式化}。将时间窗优化问题建模为马尔可夫决策过程\cite{sutton2018reinforcement}，定义状态空间（离散化的时间窗参数）、动作空间（时间窗边界调整）、奖励函数（分类准确率）和状态转移规则，构建包含 136 个状态的有限状态空间。

（2）\textbf{算法设计}。设计表格型 Q-Learning 算法，通过维护状态 - 动作价值表迭代估计最优动作价值函数，采用$\varepsilon$-贪婪策略平衡探索与利用，以 CSP-SVM 分类器的交叉验证准确率作为即时奖励\cite{watkins1992q}。

（3）\textbf{理论分析}。分析表格型 Q-Learning 的收敛性、计算复杂度和内存占用，与 DQN 进行系统性对比，明确两种方法的适用边界。

（4）\textbf{实验验证}。在 BCI Competition IV 2a 数据集上验证所提方法的有效性，从分类准确率、训练时间、推理延迟、内存占用等维度与 DQN 和固定时间窗基准进行对比。

\xmusubsection{研究意义}{Research Significance}
\label{subsec:contributions}

本研究的意义主要体现在理论贡献与实践价值两个层面：

\textbf{理论贡献}：

（1）\textbf{状态空间有限任务的算法选择准则}。本研究明确了表格型 Q-Learning 与 DQN 的适用边界：当状态空间规模小于$10^4$且状态 - 动作对可充分访问时，表格型方法在收敛性保证、样本效率和可解释性方面均优于 DQN。这一结论可推广至其他状态空间有限的 BCI 优化任务，为强化学习算法的选择提供了理论依据。

（2）\textbf{经典方法在特定场景下的再发现}。在深度学习主导的研究趋势下，本研究重新审视经典强化学习算法的价值。研究发现，在状态空间有限（136 个状态）、状态转移确定性的特定任务中，无需函数近似的表格型方法足以捕捉最优策略。这一发现挑战了"深度模型必然优于浅层模型"的固有认知，为经典方法的现代应用提供了新视角。

\textbf{实践价值}：

（1）\textbf{轻量化 BCI 设计新范式}。针对便携式 BCI 设备的资源约束，本研究提出将深度强化学习简化为表格型强化学习的设计思路。实验结果表明，表格型 Q-Learning 在保持与 DQN 相当性能（62.68\% vs 62.70\%, $p = 0.981$）的同时，实现了 3 倍训练加速、1000 倍推理延迟降低和 125 倍内存占用减少。这一结果为资源受限场景下的 BCI 系统设计提供了可行方案。

（2）\textbf{可解释性 BCI 系统的实践探索}。表格型 Q-Learning 的 Q 值表可直接可视化和分析，研究者能够追踪智能体的学习轨迹、理解决策依据并调试优化策略。在医疗康复等安全敏感应用中，可解释性是系统部署的关键前提。本研究的可解释性设计为 BCI 系统的临床转化提供了重要支撑。

\xmusection{论文结构}{Thesis Structure}
\label{sec:thesis_structure}

本论文共分为五章，结构安排如下：

\textbf{第 1 章 引言}。介绍研究背景、问题提出、研究内容和研究意义。

\textbf{第 2 章 相关工作}。综述运动想象 BCI 的特征提取、分类算法、时间窗优化策略和强化学习应用。

\textbf{第 3 章 基于表格型 Q-Learning 的时间窗优化方法}。介绍基于表格型 Q-Learning 的时间窗优化方法，包括问题形式化、算法设计和理论分析。

\textbf{第 4 章 实验与结果分析}。描述实验设置和对比方法，呈现实验结果并进行统计分析。

\textbf{第 5 章 总结与展望}。总结研究工作，讨论局限性并展望未来方向。

